\newif\ifupperformal
\upperformaltrue

\newcommand{\statusAuthorsNote}{
\ifupperformal
The formal proofs of the five results
\cref{thm:barrier-730,b:cor:subcritical,thm:upper,f:main,cor:mean-bias}
are included in the Lean development described in \cref{sec:formal}.
That section gives the compared statements, the source revisions covered
by the verification records, and the status of the strengthened
mean-bias interface. It distinguishes these formal claims from the
stronger conclusions proved only in the written argument.
\else
The formal proofs of the four results
\cref{thm:barrier-730,b:cor:subcritical,f:main,cor:mean-bias}
are included in the Lean development described in \cref{sec:formal},
which records the verification scope and the status of the strengthened
mean-bias interface. The upper-bound statement and its record are described in
\cref{sec:formal-upper}.
\fi}

\newcommand{\statusAbstract}{
\ifupperformal
We have also formalized the upper bound in continuous time.
\else
The formal treatment of the upper bound is described in
\cref{sec:formal-upper}.
\fi}

\newcommand{\statusIntroShort}{
\ifupperformal
We have formalized \cref{thm:upper}; see \cref{sec:formal-upper}.
\else
For the formal statement corresponding to \cref{thm:upper}, see
\cref{sec:formal-upper}.
\fi}

\newcommand{\statusIntroFormal}{
\ifupperformal
\Cref{thm:upper} is also formalized, as the statement described in
\cref{sec:formal-upper}, and is included in the combined five-claim
verification described in \cref{sec:formal}.
\Cref{cor:upper-discrete,cor:two-sided} are not separately formalized.
\else
\Cref{sec:formal-upper} gives the formal statement corresponding to
\cref{thm:upper}, the proof declaration, and the verification record.
\fi}

\newcommand{\statusFormalUpper}{
The proposition
\path{SKConstants.OptimizedSKTypicalStartStatement}, defined in
\path{SKConstants/SKModel.lean}, states that, for every $\beta>1$ and
every $\omega\in(0,1/4)$,
\[
 \PP_{J,X_0\sim\pi^J}
 \bigl\{d^J\bigl(X_0,\exp(n^{1-\omega})\bigr)>1/4\bigr\}
 \longrightarrow0.
\]
The declaration uses $c$ for the parameter denoted $\omega$ here, with
both $\beta$ and $c$ fixed before $n\to\infty$.

In the formal model, the variables $g_{ij}$ are independent standard
Gaussians, and the completed matrix and Hamiltonian are
\[
 W_{ij}=\frac{g_{ij}+g_{ji}}{\sqrt{2n}},\qquad
 H_g^c(x)=\frac12x^{\mathsf T}Wx,
\]
including its diagonal, with $\beta$ outside the Hamiltonian.
For $i<j$, we define $J_{ij}=\sqrt n\,W_{ij}$. These couplings are
independent standard Gaussians, and
\[
 H_g^c(x)=H_n^J(x)+\frac12\tr W,\qquad
 Z_n^c=e^{\beta\tr W/2}Z_n.
\]
The trace term is independent of $x$ and cancels in the Gibbs
normalization, as in \cref{lem:model-completion}(b), and in every energy
difference. The Gibbs probabilities and flip rates coincide with those
of \cref{sec:model}. The same holds for the semigroup and the event
defined by the total variation distance from a fixed initial spin.
The formal generator is $n(K^J-I)$, so each site has rate one, as in
\eqref{eq:common-semigroup}. Its joint law samples the completed matrix
and then a Gibbs spin. By integrating out the independent diagonal, we
obtain the joint law of \cref{def:model-joint}. The primitive representation
in \cref{u:sec:primitive} uses the same completed matrix law. The initial
spin is fixed before its transition law and total variation distance
are formed. By the complementary form of \cref{lem:model-joint}, the
displayed limit is \cref{thm:upper}. The formal sequence is also defined
at sizes zero and one, which does not affect the limit.

\ifupperformal
The theorem
\path{SKConstants.Completion.Main.optimizedSKTypicalStartStatement_holds}
proves this proposition without additional hypotheses. The original
upper-bound package at commit \path{352426b9007b6bbb31a3c8cfe6887caa0badf3fd}
was checked on 19 October 2026, using
Lean version \texttt{v4.35.0-rc2} and Mathlib at commit
\path{065356127b1dc0016f66b7283ce0ce2c4055aa55}.
A fresh build succeeded at this revision. The package's \path{Challenge.lean}
contains 47 independently stated claims and imports only Mathlib. Its
definitions explicitly include the completed disorder law, Gibbs kernel,
semigroup, total variation distance, and joint probability. The official
Comparator confirmed equality of the proved statements and these
independent statements, including their definitions. The axiom audit of
all 20,281 exported development declarations in 1,605 modules found only
\path{propext}, \path{Classical.choice}, and \path{Quot.sound}; the main
theorem has exactly these three axiom dependencies. The same Comparator
run checked all 47 claims and their proof dependencies with Lean's kernel,
NanoDa, and con-ron. All three accepted them. The upper-bound theorem
is also one of the five final statements checked in the combined package
described at the start of \cref{sec:formal}.

The formal proof follows the route of \cref{part:upper}. It constructs
Brownian motion, proves the scalar and Gaussian-query laws, establishes
the fine entropy and sharp coarse estimates, and combines coverage and
the connecting sequence with the finite-chain mixing reduction. The
steps are organized as separately proved statements and assembly
theorems, whose hypotheses are supplied by the final proof. The
parameter choice used by that proof is described in
\cref{sec:formal-scope}. The development includes Ezzeri Esa's
StochasticCalculus library and adapted OpenAI Lean modules, both under
the Apache-2.0 license. Intermediate estimates of the written proof are
checked only insofar as they enter the formal proof or are separately
stated formal results.
\else
The proof declaration and verification records are available with the
sources described in \cref{sec:formal-scope}.
\fi}

\newcommand{\statusFormalUpperScope}{
\ifupperformal
\Cref{cor:upper-discrete,cor:two-sided} are not separately formalized.
The development proves admissibility of the affine family in
\cref{u:prop:range} throughout
$0<\varepsilon_\omega<1/80$, including the choice
$\varepsilon_\omega=(1/4-\omega)/20$. The final proof of
\cref{thm:upper} instead uses the same affine family at
$\min\{1/2000,(1/4-\omega)/20\}$ and proves the target for every
$\omega\in(0,1/4)$. A separate formal result characterizes the open
range of the development's strict estimate system, and another checks
an explicit rational instance. The ancillary calculation described
after \cref{u:prop:range}, including every identity for
\cref{u:tab:budget}, is not a separately certified formal statement.
\else
The scope of the upper-bound verification is recorded with the sources
described below.
\fi}
